\subsection{Оценка зоны радиовидимости спутника по энергетике}
\label{subsec:snr}

Угол обзора спутника определяется как максимальный угол $\alpha_{\max}$
между направлением из КА в центр Земли и направлением на обслуживаемую точку на Земле.
Угол обзора может быть не зафиксирован, вместо этого он может быть 
определен требуемым качеством радиоканала. 
Качество радиоканала определяется заданным порогом отношения сигнал/шум
$\mathrm{SNR}_{\min}$. Для каждой точки
земной поверхности должно выполняться условие $\mathrm{SNR} \ge
\mathrm{SNR}_{\min}$, что напрямую ограничивает максимальную
наклонную дальность $d_{\max}$, следовательно, определяет $\alpha_{\max}$. 
Это связывает геометрию группировки с энергетическими характеристиками радиолиний.

\paragraph{Геометрия задачи и связь с наклонной дальностью}

На рисунке \ref{fig:alpha_max_geometry} представлена геометрия
задачи. КА находится на высоте $h$ над
поверхностью Земли (экваториальный радиус Земли $R_E$). Геометрическая дальность
до точки на поверхности обозначена как $d_{\max}$.

Максимальный угол $\alpha_{\max}$ находится из решения треугольника,
образованного центром Земли, спутником и точкой на поверхности, по
теореме косинусов:
\[
  \cos\alpha_{\max}(d_{\max}) = \frac{(R_{\text{Earth}} + h)^2 +
  d_{\max}^2 - R_E^2}
  {2\,(R_{\text{E}} + h)\,d_{\max}}.
\]
Таким образом, геометрия обзора напрямую связана с максимальной
наклонной дальностью $d_{\max}$, которая, в свою очередь, ограничена
энергетикой радиолинии.

\paragraph{Энергетика радиолинии}

Рассмотрим линию связи «спутник – Земля», 
в которой на космическом аппарате установлена АФАР, 
обеспечивающая электронное сканирование луча.
Для оценки энергетических характеристик канала используется 
отношение сигнал/шум ($SNR_{\min}$) на входе приемника.
Заданы параметры передающей и приемной антенн: 
мощность передатчика $P_{tx}$ (Вт), коэффициенты
усиления $G_{tx}^{(0)}$ и $G_{rx}^{(0)}$ (в дБи) в направлении
максимального излучения, шумовая температура приемника $T_n$ (К),
полоса сигнала $B$ (Гц) и несущая частота $f$ (Гц).

Расчет SNR производится по формуле:
\begin{equation}
\mathrm{SNR} = P_{tx} + G_{tx} + G_{rx}(h) - L_{\text{fs}} - P_{n}
\end{equation}

Все величины в правой части уравнения выражены в децибелах. 
В данной модели не учитываются дополнительные потери. При необходимости они могут быть учтены 
путем добавления соответствующего запаса к пороговому значению \(\mathrm{SNR}_{\min}\).

Согласно формуле Фрииса~\cite{friis1946}, потери в свободном
пространстве на наклонной дальности $d$ составляют (в дБ):
\begin{equation}\label{eq:fs_loss}
L_{\text{fs}}(d) = 20\log_{10}\!\left(\frac{4\pi d}{\lambda}\right),
\end{equation}
где $\lambda = c/f$ -- длина волны, $c$ -- скорость света.

Мощность шума на входе приемной антенны (в дБВт):
\begin{equation}
P_{n} = 10\log_{10}(k\,T_n\,B),
\end{equation}
где $k$ -- постоянная Больцмана.

При отклонении луча от нормали антенны усиление снижается за 
счёт потерь на сканирование. Для оценки границы зоны обслуживания 
используются эффективные значения коэффициентов усиления на краю 
зоны радиовидимости \eqref{eq:gain_edge}. Предельно допустимый угол 
сканирования АФАР $\psi_{\max}$ определяется исходя из ограничений, 
связанных с падением эффективности решётки при отклонении луча, 
как показано в классической работе~\cite{balanis2016}. 
Использование такого порогового значения позволяет рассматривать 
зависимость угла обзора от высоты как квадратичную функцию, 
что делает возможным аналитическое решение задачи. 
Это критично для скорости работы оптимизатора орбитальной группировки.

\begin{equation}\label{eq:gain_edge}
G^{\text{eff}} = G^{(0)} + 10\log_{10}(\cos \psi^{\max}).
\end{equation}

На границе зоны радиовидимости выполняется равенство $\mathrm{SNR} =
\mathrm{SNR}_{\min}$. С учетом потерь в свободном пространстве \ref{eq:fs_loss} и
эффективных усилений \ref{eq:gain_edge} получаем:
\begin{equation}
\mathrm{SNR}_{\min} = P_{tx} + G_{tx}^{\text{eff}} + G_{rx}^{\text{eff}}
- 20\log_{10}\!\left(\frac{4\pi d_{\max}}{\lambda}\right) - P_{n}.
\end{equation}

Тогда максимальная наклонная дальность выражается следующим образом:
\begin{equation}
 d_{\max}(h) = \frac{\lambda}{4\pi} \cdot 10^{\frac{M -
 \mathrm{SNR}_{\min}}{20}},
\end{equation}
где
\begin{equation}
M = P_{tx} + G_{tx}^{\text{eff}} + G_{rx}^{\text{eff}} - P_{n}
\end{equation}


\paragraph{Влияние на выбор орбитального построения}

Таким образом, зависимость $\alpha_{\max}(h)$ учитывает и 
энергетику радиолинии и геометрические ограничения.
Число спутников, необходимое для глобального непрерывного покрытия,
обратно пропорционально величине угла обзора: чем шире зона радиовидимости
каждого спутника, тем меньшим количеством спутников можно покрыть всю
земную поверхность.
При понижении высоты КА уменьшаются потери в свободном
пространстве, что позволяет увеличить максимальную наклонную
дальность $d_{\max}$ при заданном $\mathrm{SNR}_{\min}$.
Соответственно, увеличивается и угол обзора $\alpha_{\max}$ (вплоть
до технического предела $\psi_{\max}$). 
Однако понижение высоты при фиксированном угле обзора требует
большего количества КА для обеспечения непрерывного покрытия. 

Следовательно, оптимизатор орбитальной группировки балансирует
между увеличением зоны радиовидимости (за счет
роста $\alpha_{\max}$ при снижении высоты) и увеличением числа
спутников, необходимого для полного покрытия Земли при понижении высоты. Это
позволяет находить высоту спутниковой группировки, при которой требуемое
количество КА в группировке минимально при соблюдении заданного
качества связи $\mathrm{SNR}_{\min}$.


